\section{The arithmetic cutoff comparison}
This section distinguishes the zero-side calculation from
the estimates required for the arithmetic prolate image. The full
cutoff limit requires the frequency-tightness estimate stated below.
For a fixed primitive Dirichlet character, let $Q_\Lambda$ be the
orthogonal projection onto its selected theta images, normalized by
their full Gram matrix, and let $S_L=1_{[-L,L]}$, $L=\log\Lambda$.
The compressed defect $D_L=S_L-S_LQ_\Lambda S_L$ is positive.
To identify its trace on every compact smooth logarithmic test,
one needs local convergence of its Fourier-diagonal measure and a
uniform weighted tail estimate. Positivity and a finite-packet lower
bound alone do not provide the latter estimate.

The unconditional one-sided closure is known: Burnol's
Theorem 3.2 identifies it with a Hardy space multiplied by the
Blaschke product of zeros in the right half-plane
\citep[pp.~14--15]{burnol2000}. This theorem concerns a one-sided
closure, not a finite prolate selection. The principal sector also
requires its mean and evaluation constraints and the two polar
coordinates. The following sections prove these comparisons for
this same cutoff, including its transition eigenfunctions.
